\section{Statistical theory}
\label{sec:theory}

\subsection{Risks and what held-out labels can measure}

Two risks are relevant. The first is measured against the estimand; the second against the
per-comparison target $\mu_X=\E\{\phi(Y)\mid X\}$:
\[
  R(\hmu)=\E\norm{\hmu(Z)-\muH(Z)}^2,\qquad
  R_X(\hmu)=\E\norm{\hmu(Z)-\mu_X}^2=R(\hmu)+B_\zeta .
\]
The expectation is over the evaluation distribution $Z\sim P_Z$ (density $q$) and over the training
sample. The equality holds because $\E\{\mu_X-\muH(Z)\mid Z\}=0$. With the embeddings of
Section~\ref{sec:setup}, $R$ is the Brier/MSE risk for binary labels, the ranked probability score (a
discrete CRPS) for ordinal labels, and a conditional MMD$^2$ for kernel embeddings
\citep{gneiting2007strictly,muandet2017kernel}. \citet{li2026personalizing} integrate against
Lebesgue measure on $[0,1]^d$. When $q$ is bounded above and below, the two risks are equivalent up to
constants, but their optimal designs differ (Section~\ref{sec:acquisition}).

\begin{lemma}[Single held-out labels suffice for comparisons]
\label{lem:heldout}
For a test comparison independent of the training data,
$\E\norm{\phi(Y)-\hmu(Z)}^2=R_X(\hmu)+\E\tr\Cov\{\phi(Y)\mid X\}$.
\end{lemma}

\begin{proof}
Write $\phi(Y)-\hmu=\{\phi(Y)-\mu_X\}+\{\mu_X-\hmu\}$. Conditional on $X$ and the training data, the
cross term has mean zero.
\end{proof}

The last term does not depend on the method. Differences in held-out Brier score or RPS on the
common test set $\mathcal T$ are therefore unbiased for differences in $R_X$, even with one human
label per comparison. Metrics that need $\mu_X$ itself require multi-annotator comparisons.

\subsection{When does the full judge distribution help?}
\label{sec:scalar-vs-dist}

\begin{proposition}[Scalar versus distributional localization]
\label{prop:scalar-vs-dist}
Let $Z_s=\zeta_s(X)$ and $Z_d=\zeta_d(X)$, with $Z_s$ a measurable function of $Z_d$ (for example,
$Z_s=(p_J,\xi)$ and $Z_d=(p_J,t_J,o_J,\xi)$). Let $\mu_s^{*},\mu_d^{*}$ be the corresponding
estimands, and $\hmu_s,\hmu_d$ estimators fitted on data independent of the evaluation point. Then
\[
  R_X(\hmu_s)-R_X(\hmu_d) \;=\; G \;-\;\bigl\{R(\hmu_d)-R(\hmu_s)\bigr\},
  \qquad G=\E\norm{\mu_d^{*}(Z_d)-\mu_s^{*}(Z_s)}^2\ \ge 0,
\]
where each $R(\cdot)$ is measured against its own estimand. Moreover, $G=0$ if and only if
$\E\{\phi(Y)\mid Z_d\}=\E\{\phi(Y)\mid Z_s\}$ almost surely.
\end{proposition}

\begin{proof}
$R_X(\hmu_\bullet)=R(\hmu_\bullet)+\E\norm{\mu_\bullet^{*}(Z_\bullet)-\mu_X}^2$. The $\sigma$-fields are
nested, so $\mu_s^{*}(Z_s)=\E\{\mu_d^{*}(Z_d)\mid Z_s\}$, and orthogonality gives
$\E\norm{\mu_s^{*}(Z_s)-\mu_X}^2=\E\norm{\mu_d^{*}(Z_d)-\mu_X}^2+G$. Subtract the two expressions.
\end{proof}

\begin{remark}[Status and use]
The identity is standard conditional-expectation orthogonality. For binary labels, $G$ is the
numerator of the nonparametric variable-importance measure of the extra judge features, in the
sense of \citet{williamson2021nonparametric,williamson2023general}. The corresponding monotonicity
of oracle-calibrated Brier risk for nested judge panels is Proposition~1 of \citet{li2026calibrate}.
We claim no novelty for it. Its role here is operational. First, it names the estimand, $G$, that
decides whether the distributional channel carries information. Second, it shows that comparing
two \emph{fitted} predictors estimates $G+R(\hmu_s)-R(\hmu_d)$, not $G$. Testing $G=0$ therefore needs
variable-importance inference (experiment E2), not a comparison of fitted scores. Rate upper bounds
for $R(\hmu_s)$ and $R(\hmu_d)$ also do not determine the sign or order of their difference, so the
crossover budget is an empirical quantity.
\end{remark}

\subsection{Finite judge samples: oracle information loss}
\label{sec:judge-budget}

If $\QJ$ is elicited from $M$ stochastic calls (protocol J2), we observe a noisy summary $\hat Z_M$
instead of $Z_\infty$.

\begin{proposition}[Judge budget; oracle information loss]
\label{prop:judge-budget}
Suppose:
\begin{itemize}
  \item the judge's sampling noise is independent of the comparison and the human label given
        $Z_\infty$, that is, $\hat Z_M\perp(X,Y)\mid Z_\infty$;
  \item $\E\norm{\hat Z_M-Z_\infty}^2\le C_\psi/M$;
  \item $\mu_\infty^{*}(z)=\E\{\phi(Y)\mid Z_\infty=z\}$ extends to a function on a common domain
        containing the supports of both $Z_\infty$ and $\hat Z_M$, and satisfies
        $\norm{\mu_\infty^{*}(z)-\mu_\infty^{*}(z')}\le\theta_1^{*}\norm{z-z'}^{\theta_2^{*}}$ there.
\end{itemize}
Then the best predictor measurable with respect to $\hat Z_M$ loses at most
\[
  \E\norm{\E\{\phi(Y)\mid\hat Z_M\}-\mu_X}^2-\E\norm{\mu_\infty^{*}(Z_\infty)-\mu_X}^2
  \;\le\;(\theta_1^{*})^2\,(C_\psi/M)^{\theta_2^{*}} .
\]
\end{proposition}

\begin{proof}
Conditional independence gives $\hat Z_M\perp Y\mid X$, so
$\E\{\phi(Y)\mid\hat Z_M\}=\E\{\mu_X\mid\hat Z_M\}$. This is the $L^2$ projection of $\mu_X$ onto
functions of $\hat Z_M$, so its error is at most that of $\mu_\infty^{*}(\hat Z_M)$. Expand
$\mu_\infty^{*}(\hat Z_M)-\mu_X=\{\mu_\infty^{*}(\hat Z_M)-\mu_\infty^{*}(Z_\infty)\}+\{\mu_\infty^{*}(Z_\infty)-\mu_X\}$.
The cross term vanishes because $\E\{\mu_X\mid Z_\infty,\hat Z_M\}=\mu_\infty^{*}(Z_\infty)$. Finally,
use the H\"older bound and Jensen's inequality:
$\E\norm{\hat Z_M-Z_\infty}^{2\theta_2^{*}}\le(C_\psi/M)^{\theta_2^{*}}$.
\end{proof}

This result has three limits.
\begin{itemize}
  \item It bounds \emph{oracle} information loss only. It does not show that the kernel estimator
        run on $\hat Z_M$ attains that oracle risk plus $r_n$: noise also changes the localization
        geometry and the offset. E5 is therefore exploratory.
  \item Conditional independence holds when $\hat Z_M$ is computed from $M$ i.i.d.\ draws whose law
        is determined by $Z_\infty$. This is true if $Z_\infty$ contains the full judge probability
        vector, but not automatically for a compressed summary.
  \item $M\gtrsim r_n^{-1/\theta_2^{*}}$ makes the loss comparable to an estimation rate $r_n$;
        negligibility needs $M\gg r_n^{-1/\theta_2^{*}}$.
\end{itemize}

\subsection{Target results: rates, lower bound, adaptation}

\begin{assumption}[Regularity]
\label{ass:reg}
The continuous part of $Z$ ranges over $\cZ=[0,1]^d$ (after stratifying on discrete features).
\begin{enumerate}[label=(\alph*)]
  \item $\norm{\muH(z)-\muH(z')}_{\cH}\le\theta_1^{*}\norm{z-z'}^{\theta_2^{*}}$ with
        $0<\theta_2^{*}\le1$.
  \item $\norm{\phi}\le1$. Write $\sigma^2(z)=\E\{\norm{\phi(Y)-\muH(z)}^2\mid Z=z\}$; $\sigma(\cdot)$ is
        H\"older.
  \item \emph{Nondegeneracy:} $\sigma(z)\ge\sigma_{\min}>0$. For binary labels this holds when
        $p_H^{*}\in[\kappa,1-\kappa]$, with $\sigma_{\min}^2=\kappa(1-\kappa)$.
  \item \emph{Design and occupancy:} revealed points have a design density $p_Z$ with
        $c_p\le p_Z\le C_p$ on $\cZ$; the evaluation density $q$ satisfies $c_q\le q\le C_q$; and
        the bandwidth satisfies $nh^d\ge C\log n$.
\end{enumerate}
\end{assumption}

Conditions (c) and (d) are a simple sufficient replacement for the low-noise-region condition in
Theorem~1 of \citet{li2026personalizing}. Without such a condition the rate below is false.
\emph{Counterexample:} take $d=1$, $\mJ(z)=z$, $Y\mid Z=z\sim\Ber(n^{-3})$. The prescribed bandwidth is
then $h\asymp n^{-4/3}$, so almost all windows are empty. On an empty window the estimator returns
$\mJ(z)=z$, and its risk tends to $1/3$, while the displayed rate is $O(n^{-8/3})$.

For a fixed smoothing parameter $\theta$, let $\mathcal F_\theta(\theta^{*},\gamma,\sigma_0)$ be the laws
satisfying Assumption~\ref{ass:reg}(a)--(c) with $|\cZ|^{-1}\int\sigma\le\sigma_0$ and
$\delta_{\theta,z}\in\mathcal L_z(\gamma)$ for every $z$. Here $\delta\in\mathcal L_z(\gamma)$ means
$\sup_{z'}\norm{\delta(z')-\delta(z)}/\norm{z'-z}^{\gamma_2}\le\gamma_1$. Write
$\bar\sigma=|\cZ|^{-1}\int\sigma$.

\begin{target}[Upper bound, fixed $\theta$]
\label{tr:upper}
Under Assumption~\ref{ass:reg}, with $\gamma=\gamma(\theta)$ and
$h\asymp\min\{(\bar\sigma/\gamma_1)^{2/(2\gamma_2+d)}n^{-1/(2\gamma_2+d)},1\}$ subject to
$nh^d\ge C\log n$,
\[
  R(\hmu_{\theta,h})\;\lesssim\;\gamma_1^{\frac{2d}{2\gamma_2+d}}\,\bar\sigma^{\frac{4\gamma_2}{2\gamma_2+d}}\,
  n^{-\frac{2\gamma_2}{2\gamma_2+d}}\;+\;\frac{\bar\sigma^2}{n}.
\]
\end{target}

\begin{target}[Minimax lower bound for bounded preference probabilities]
\label{tr:lower}
Consider binary labels, and a \emph{fixed} judge (not part of the supremum) with interior slack,
$\mJ(z)\in[2\kappa,1-2\kappa]$. Let $\mathcal F_\theta^{\kappa}$ be the subclass of
$\mathcal F_\theta(\theta^{*},\gamma,\cdot)$ with $p_H^{*}\in[\kappa,1-\kappa]$, and assume
$\gamma_1\le\kappa$. Take $0\le\theta_1\le C\theta_1^{*}$ and $\theta_2\ge\theta_2^{*}$. Then for
$n\ge n_0(\kappa)$, the infimum over designs and over estimators with access to $\mJ$ satisfies
\[
  \inf_{p_Z,\hat p}\;\sup_{\mathcal F_\theta^{\kappa}}\;R(\hat p)\;\gtrsim_\kappa\;
  \gamma_1^{\frac{2d}{2\gamma_2+d}}\,n^{-\frac{2\gamma_2}{2\gamma_2+d}}+\frac1n .
\]
\end{target}

The slack is necessary. With $\mJ(z)=z$ on $[0,1]$ and $\gamma_1=0$, the class forces
$p_H^{*}(z)=z+c$, and the endpoint constraints force $c=0$. The class is then a single point with
minimax risk zero, although $\bar\sigma=\pi/8$. For binary labels, $\bar\sigma$ is not a free
parameter, since it is determined by $p_H^{*}$, which is why the bound is stated without it. The
proof must build the packing inside the slack: bump amplitudes $\le\gamma_1h^{\gamma_2}\le\kappa$, and
the parametric part from constant shifts of size $n^{-1/2}\le\kappa$.

\begin{target}[Adaptation and no harm]
\label{tr:oracle}
$R(\hmu_{\hat\theta,\hat h})\le C_1\min_{(\theta,h)\in\Theta\times\mathbb H}R(\hmu_{\theta,h})
+C_2\sigma_{\max}^2(\log n)^2/n$.
\end{target}

\emph{Scope of no harm.} The comparison is with the human-only kernel estimator \emph{using the same
localization variable $Z$ and the same revealed sample}. Selecting $\theta_1=0$ switches off the
offset. It cannot restore information that $Z$ discards. If the judge is corrupted, $Z$ is
corrupted too, because $Z$ contains judge summaries.

\paragraph{When is the judge helpful?}
As in Examples~1--2 of \citet{li2026personalizing}, borrowing helps when the discrepancy has a
smaller H\"older constant than the target ($\gamma_1(\theta)\ll\theta_1^{*}$; for example, a nearly
calibrated judge) or a higher smoothness order ($\gamma_2(\theta)>\theta_2^{*}$).

\begin{remark}[Which theory is routine and which carries the paper]
\label{rem:novelty-theory}
\citet{li2026personalizing} already cover binary classification (their Section~5.2), so for a
binary human label their scalar predictor already estimates a Bernoulli distribution. Targets
\ref{tr:upper} and \ref{tr:oracle} should be routine coordinatewise extensions. Enlarging $K$ or
using a kernel embedding brings no new variance phenomenon, since $\norm{\phi}\le1$ makes averaging
dimension-free. The theoretical weight of the paper therefore rests on three results:
\begin{itemize}
  \item Target~\ref{tr:lower}: minimax lower bounds for \emph{bounded} preference targets with a
        fixed judge, where boundary constraints interact with the judge;
  \item Target~\ref{tr:design}: finite-pool adaptive design under the population risk;
  \item Target~\ref{tr:e2e}: the end-to-end minimax rate for judge-assisted preference
        optimization, which combines both with the bridge of Section~\ref{sec:post-training}.
\end{itemize}
The propositions in Sections~\ref{sec:scalar-vs-dist}--\ref{sec:judge-budget} and~\ref{sec:post-training}
are elementary and serve as tools.
\end{remark}
